\chapter{Silica, Wood Dust and Respiratory Hazard Risk Assessment}
\label{chap:Silica, Wood Dust and Respiratory Hazard Risk Assessment}

%---------------------------- SECTION ----------------------------
\section{Why this assessment matters in construction}

Occupational lung disease is the most significant source of work-related mortality in UK construction after asbestos-related disease. Silicosis, chronic obstructive pulmonary disease (COPD), and occupational lung cancer caused by respirable crystalline silica (RCS) exposure, together with occupational asthma and sinonasal cancer caused by wood dust exposure, kill and disable construction workers on a scale that is almost entirely preventable. The HSE's Breathe Freely campaign has quantified the impact: approximately 500 construction workers die each year from COPD alone caused by occupational dust exposure; thousands more develop silicosis, an irreversible and progressive lung disease caused by the accumulation of silica dust in lung tissue that is entirely distinct from and additional to COPD.

Respirable crystalline silica is present in very high concentrations in the materials that construction operatives cut, grind, and drill on a daily basis. Sandstone and granite typically contain over 70 per cent crystalline silica by weight; brick contains between 20 and 30 per cent; concrete contains between 15 and 25 per cent; and calcium silicate blocks, mortar, and many ceramic tiles contain significant crystalline silica fractions. The critical feature of RCS is that the particles that cause lung disease (particles smaller than 10 micrometres in aerodynamic diameter, defined as respirable) are invisible to the naked eye. An operative cutting concrete block with an angle grinder does not see the cloud of respirable silica particles that surrounds them; they see the visible coarser dust, which settles quickly, and experience the visible dust as a nuisance rather than a hazard.

The UK workplace exposure limit for respirable crystalline silica is 0.1\,mg/m$^3$ as an eight-hour time-weighted average. Dry cutting of concrete, brick, or sandstone with an angle grinder routinely generates RCS air concentrations of between 1 and 20\,mg/m$^3$ at the operative's breathing zone --- ten to two hundred times the WEL. The legal maximum is not a generous limit; it represents the level below which the HSE's scientific advisors consider that the risk of serious harm is adequately controlled. Exposure at ten times the WEL, sustained over a working career, produces silicosis with high certainty.

Wood dust from hardwoods is classified by the International Agency for Research on Cancer as a Group 1 carcinogen (confirmed cause of cancer in humans), specifically associated with sinonasal cancer and adenocarcinoma of the nasal cavity. The WEL for hardwood dust is 1\,mg/m$^3$; for softwood dust it is 5\,mg/m$^3$. Uncontrolled routing, sanding, and cutting of hardwood timbers and MDF (which contains wood dust as well as formaldehyde binder resin) generates dust levels that frequently exceed the WEL.

\barquote{``The dust that kills is the dust you cannot see. The visible cloud settles; the respirable fraction that stays airborne is the fraction that enters and stays in the lungs.''}

\example{Site Reality}{A groundworks and landscaping contractor employs operatives on a two-year contract laying granite sett paving across a town centre regeneration scheme. The setts are cut to size using dry angle grinders. No respiratory controls are specified in the COSHH assessment beyond the statement ``dust masks to be worn where dust is generated''. The ``dust masks'' supplied are disposable FFP2 dust masks; FFP2 provides approximately 10-fold protection --- insufficient for RCS exposures that are 10--200 times the WEL generated by dry angle grinding of granite. Five years after the contract ends, one operative is diagnosed with silicosis at the age of 39. He has progressive pulmonary fibrosis; the silicosis has no treatment that can reverse or halt progression. The COSHH assessment failed to identify the exposure concentration, failed to apply the WEL, failed to consider water suppression and on-tool extraction, and failed to specify FFP3 minimum RPE for an inadequate control that would have been insufficient regardless.}

%---------------------------- SECTION ----------------------------
\section{Legal and professional framework}

The Control of Substances Hazardous to Health Regulations 2002 provide the primary framework for respiratory hazard management in construction, requiring that exposure to hazardous substances including RCS and wood dust is assessed and controlled to below the relevant WEL where a WEL has been assigned. EH40 (Workplace Exposure Limits) assigns WELs for respirable crystalline silica (0.1\,mg/m$^3$ TWA), hardwood dust (1\,mg/m$^3$ TWA), and softwood dust (5\,mg/m$^3$ TWA). The CDM 2015 pre-construction information duties require designers to identify materials in the existing structure that may generate hazardous dust.

Health surveillance is required under COSHH where there is a reasonable likelihood of an identifiable disease or adverse health effect arising from work with hazardous substances: for RCS and hardwood dust, health surveillance is specifically indicated because both cause progressive and irreversible lung disease with a long latency. The HSE Breathe Freely campaign, developed in partnership with construction industry bodies, provides practical tools including task-based exposure assessments and control solutions for construction dust.

\paragraph{Key frameworks relevant to this chapter include: COSHH Regulations 2002; EH40; RIDDOR 2013; HSE Breathe Freely guidance; CDM 2015.}

\begin{itemize}
  \bitem{COSHH Regulations 2002 and EH40}{Require that exposure to RCS and wood dust is maintained below the relevant WEL; require health surveillance for workers exposed to substances that cause progressive occupational disease; require that control measures are adequately maintained; require biological monitoring and/or exposure monitoring where necessary to verify control effectiveness.}
  \bitem{HSE Breathe Freely Campaign and Task-Based Assessments}{Provide the construction industry with task-specific dust exposure data and a toolbox of control solutions; the task-based assessment approach allows the COSHH assessor to estimate RCS and wood dust exposure levels from known task and control combinations without requiring individual air monitoring for routine activities.}
  \bitem{RIDDOR 2013}{Requires reporting to the HSE of all cases of occupational asthma (as a prescribed disease) and silicosis (as a prescribed disease); health surveillance findings that identify a case of occupational lung disease must trigger RIDDOR reporting, investigation, and reassessment of the COSHH controls.}
  \bitem{CDM 2015}{Requires designers to consider the dust hazard implications of their specification choices: specifying silica-containing materials where substitutes are available without justification for the choice represents a failure of the CDM designer duty to eliminate risks through design.}
\end{itemize}

%---------------------------- SECTION ----------------------------
\section{Conducting the assessment using the five-step method}

\subsection{Step 1 --- Identify hazards}
\paragraph{Respiratory dust hazards in construction include: respirable crystalline silica from dry cutting, grinding, drilling, and scabbling of concrete, brick, block, stone, and tile; RCS from dry sweeping of concrete or masonry dust on floors; wood dust from routing, turning, sanding, and cutting of hardwood and softwood timber, MDF, and chipboard; nuisance dust from insulation materials including mineral wool, glass fibre, and Rockwool; mixed construction dust from demolition operations generating silica-containing material; and organic dust from agricultural soils encountered in groundworks. Each hazard must be assessed by task: the exposure generated from a specific combination of material, tool, and control measure is the unit of analysis, not a general assessment of ``dusty work''.}

\subsection{Step 2 --- Decide who or what is affected}
\paragraph{Persons at risk from construction respiratory hazards include: all operatives directly engaged in dust-generating activities; operatives in adjacent areas who are exposed to dust generated by another operative's activity; operative supervisors and site managers in proximity to dusty activities; and members of the public adjacent to construction sites where dust is not adequately suppressed. Particularly vulnerable individuals include: operatives with existing respiratory conditions (asthma, COPD) whose conditions may be aggravated by dust exposure; younger workers and apprentices who will accumulate greater lifetime exposure; and any worker already sensitised to a dust-associated sensitiser (wood dust sensitisation; isocyanate sensitisation).}

\subsection{Step 3 --- Evaluate likelihood and consequence}
\paragraph{The consequence of uncontrolled long-term RCS or hardwood dust exposure is irreversible progressive lung disease (silicosis, COPD, cancer), which is career-ending and potentially fatal. The probability of exposure at levels above the WEL for dry cutting activities without controls is effectively certain (probability 5 on a 5×5 scale) given the documented exposure data. Inherent risk for uncontrolled dry cutting must be rated Very High (25/25). With full implementation of on-tool water suppression, on-tool extraction (OTE), and FFP3 RPE, the residual risk can be reduced to Low (4--6/25). The assessment must reference the task-based exposure data available from the HSE Breathe Freely resources to support the risk rating.}

\subsection{Step 4 --- Select and implement controls}
\paragraph{Controls for RCS and wood dust must follow the COSHH hierarchy. Elimination: avoid producing dust by specifying pre-cut materials delivered to size; avoid dry operations by designing for wet-cut construction details. Substitution: use low-silica alternatives (calcium silicate blocks have lower RCS content than dense aggregate blocks; timber alternatives to stone for aesthetic features). Engineering controls: on-tool water suppression systems on all disc cutters and grinders used on silica-containing materials; on-tool extraction (OTE) with HEPA-filtered vacuums for enclosed activities; vacuum-equipped sanding tools for wood dust; prohibition of dry sweeping (replacing with vacuum collection or wet methods). PPE: where engineering controls are insufficient to achieve compliance with the WEL (high-frequency tasks, high silica-content materials, restricted spaces preventing OTE use), RPE must be specified at minimum FFP3 for RCS (FFP2 provides inadequate protection); health surveillance must be arranged for all operatives with regular exposure.}

\subsection{Step 5 --- Record and review}
\paragraph{The COSHH assessment for dust must record the specific tasks, materials, tools, and controls assessed; must reference the task-based exposure data used to estimate the exposure level; must specify the controls required for each task; and must identify the health surveillance arrangements. The assessment must be reviewed when new dust-generating materials or tasks are introduced, when health surveillance identifies an adverse finding, and at minimum annually.}

\example{Five-Step Application}{A stone restoration contractor is cutting and grinding Portland stone for a heritage building repair project. Step 1 identifies: RCS from dry angle grinding of Portland stone (65\% crystalline silica); RCS from dry disc cutting of stone setts. Step 2 identifies: two stone masons; scaffolder in proximity during scaffold inspection. Step 3 rates inherent dry grinding RCS risk as Very High (25/25); estimated exposure without controls 15\,mg/m$^3$ vs WEL 0.1\,mg/m$^3$. Step 4 implements: water suppression system fitted to all angle grinders; water-cooled disc cutter specified; FFP3 half-mask respirator worn for all grinding and cutting; dust-free area designated for stone masons during grinding; health surveillance (lung function testing) arranged at six-month intervals for all exposed operatives. Residual risk rated Low (6/25) with all controls in place. Step 5: health surveillance records maintained; OTE vacuum filter inspected monthly.}

%---------------------------- SECTION ----------------------------
\section{Applying the hierarchy of controls}

\paragraph{The hierarchy for respiratory dust control in construction is well established and supported by extensive evidence: on-tool water suppression and extraction are engineering controls of demonstrated effectiveness; RPE is the final layer, not the primary control.}

\begin{itemize}
  \bitem{Elimination}{Specify pre-cut materials delivered to exact dimensions to eliminate on-site cutting; design construction details that can be achieved without silica-containing materials; specify factory-applied finishes to avoid on-site sanding and grinding.}
  \bitem{Substitution}{Specify lower-silica materials where alternatives exist: calcium silicate block rather than dense aggregate block; timber cladding rather than masonry; resin-bonded rather than silica-bonded abrasive cutting discs; MDF with lower formaldehyde binder content.}
  \bitem{Engineering controls}{On-tool water suppression on all disc cutters and angle grinders used on silica-containing materials; on-tool extraction (OTE) vacuum with HEPA filtration for enclosed grinding, routing, and sanding; exhaust ventilation at cutting benches; prohibition of dry sweeping; enclosed cutting booths for bench-based operations.}
  \bitem{Administrative controls}{Restrict duration of high-exposure tasks; rotate operatives between high-exposure and low-exposure tasks to reduce individual daily exposure; exclude operatives from adjacent areas during high-dust activities; implement health surveillance programme; train all operatives in dust hazards, the invisible nature of the respirable fraction, and the correct use of controls; prohibit dry sweeping across the site.}
  \bitem{PPE}{FFP3 minimum disposable or half-mask respirator for all RCS-generating activities (FFP2 is insufficient); P3-rated half-mask respirator for regular wood dust exposure; powered air purifying respirator (PAPR) with P3 filter for high-duration or high-intensity exposures; fit testing for all non-disposable RPE; maintenance and replacement schedule for all respiratory protective equipment.}
\end{itemize}

%---------------------------- SECTION ----------------------------
\section{Cadence of assessment and review triggers}

\paragraph{The respiratory hazard assessment must be reviewed at each change in materials, tools, or working methods, and the health surveillance programme provides the ongoing monitoring mechanism that alerts the assessor to any failure of the controls to maintain exposure below levels causing harm.}

\begin{itemize}
  \bitem{Before each new dust-generating work phase}{When a new material or activity is introduced (new stone type, new timber species, new cutting or grinding method), the COSHH assessment for that activity must be reviewed or created before the activity commences.}
  \bitem{When health surveillance identifies an adverse finding}{Any health surveillance finding indicating declining lung function, new occupational sensitisation, or onset of respiratory symptoms must trigger an immediate COSHH reassessment, a review of control effectiveness, and consideration of whether the affected operative can continue their current duties.}
  \bitem{When new evidence or guidance is published}{The HSE Breathe Freely task-based exposure database is updated periodically with new task/control combinations; any revision of the WEL for RCS or wood dust must trigger a review of all relevant COSHH assessments.}
  \bitem{Following an air monitoring result above the WEL}{If air monitoring reveals exposure above the WEL for any dust hazard, work must be suspended until controls are reviewed and improved; the assessment must be updated before the activity resumes.}
  \bitem{At minimum annually}{All dust COSHH assessments for ongoing activities must be formally reviewed at least annually to confirm that activities, materials, and controls remain as described.}
\end{itemize}

%---------------------------- SECTION ----------------------------
\section{Common failure modes}

\begin{itemize}
  \bitem{FFP2 specified for RCS activities}{FFP2 respirators are commonly specified for all ``dusty'' construction activities; they provide a nominal protection factor of 10, which is insufficient for RCS exposures generated by dry cutting (typically 10--200 times the WEL). FFP3 must be specified for all RCS activities; the difference in protection factor is ten-fold.}
  \bitem{On-tool water suppression purchased but not used}{Water suppression attachments are available on site and specified in the COSHH assessment but are not connected or used in practice because they slow down the work, create wet working conditions, or require more frequent disc changes. The control must be enforced, not offered as optional.}
  \bitem{Dry sweeping of silica dust}{Dry brooms are used to sweep up silica-containing dust from cutting and grinding activities, resuspending settled respirable particles that have accumulated on floor and work surfaces. Prohibition of dry sweeping must be enforced as a site rule.}
  \bitem{Health surveillance not arranged}{The COSHH assessment acknowledges the need for health surveillance for regular RCS and hardwood dust exposure but health surveillance is not arranged because it is perceived as expensive or logistically difficult for itinerant construction workers.}
  \bitem{Task-based assessment not conducted}{The COSHH assessment uses a generic ``dust'' entry rather than a task-specific assessment for each dust-generating activity; the specific material, tool, frequency, and control combination is not assessed and the exposure level is not estimated.}
  \bitem{RPE not fit-tested}{Half-mask respirators are supplied to operatives without fit testing; operatives with facial hair, missing teeth, or unusual facial geometry are wearing respirators that do not seal adequately and are providing minimal protection.}
\end{itemize}

\example{Failure Pattern}{A joinery contractor is manufacturing bespoke oak window frames in a site-based workshop. The workshop contains four workbenches with router tables, a thickness planer, and a disk sander. No LEV is installed; air extraction consists of two open windows at the rear of the workshop. Hardwood dust accumulates on every surface and is swept with a dry brush at the end of each day. The four joiners are assessed at the induction as requiring ``dust masks when sanding''. Three years into the contract, two joiners report nasal symptoms; one is subsequently diagnosed with early sinonasal pathology. The investigation finds sustained hardwood dust exposure well above the 1\,mg/m$^3$ WEL across the preceding three years; no LEV, no health surveillance, inadequate RPE specification, and dry sweeping contributing to persistent resuspension of settled dust.}

%---------------------------- CHAPTER SUMMARY ----------------------------
\section{Chapter Summary}

\begin{table}[H]
\centering
\begin{tabularx}{\textwidth}{>{\raggedright\arraybackslash}p{3.2cm}X}
\toprule
\textbf{Assessment Element} & \textbf{Operational Requirement} \\
\midrule
Legal/Professional Basis & COSHH Regulations 2002; EH40 (RCS WEL 0.1\,mg/m$^3$; hardwood dust WEL 1\,mg/m$^3$); HSE Breathe Freely guidance. Task-based exposure assessment required; health surveillance indicated for regular exposure. \\
Method & Apply the full five-step process and document inherent/residual risk. \\
Controls & On-tool water suppression and OTE as primary engineering controls; FFP3 RPE for RCS; prohibition of dry sweeping; health surveillance programme for regularly exposed workers. \\
Cadence & Before each new dust-generating phase; when health surveillance identifies adverse findings; following air monitoring results above WEL; minimum annually. \\
Assurance & Use audit, incident learning, and governance reporting to verify effectiveness. \\
\bottomrule
\end{tabularx}
\end{table}

\begin{itemize}
  \bitem{Key takeaway 1}{Dry cutting of concrete, brick, block, and stone generates RCS at concentrations 10 to 200 times the WEL; on-tool water suppression is a non-negotiable engineering control, not a desirable enhancement, for all dry cutting of silica-containing materials.}
  \bitem{Key takeaway 2}{FFP2 respirators are insufficient for RCS exposure; FFP3 is the minimum specification; all non-disposable respirators must be fit-tested to confirm an adequate face-seal for the individual wearer.}
  \bitem{Key takeaway 3}{Hardwood dust is an IARC Group 1 carcinogen; health surveillance including regular lung function assessment is indicated for all construction workers with regular hardwood dust exposure; health surveillance is the early warning system that identifies control failures before irreversible disease develops.}
  \bitem{Key takeaway 4}{Dry sweeping of construction dust must be prohibited across the site and replaced with vacuum collection or wet methods; dry sweeping resuspends settled respirable particles and extends exposure beyond the cutting and grinding activity itself.}
\end{itemize}

\barquote{In Chapter~\ref{chap:Electrical Safety Risk Assessment}, we turn from this assessment to the next domain in the Prudentia framework.}
